\section{Moments of symmetric cokernels}

Fix a prime $p$ and an integer $k\ge1$, and put $R=\Z/p^k$. Let $G$ be a finite abelian $p$-group with $p^kG=0$. By the classification of finite abelian groups in Volume~3,
\[
G\cong\Z/p^{\lambda_1}\oplus\dots\oplus\Z/p^{\lambda_r},\qquad1\le\lambda_i\le k,
\]
with generators $g_1,\dots,g_r$; $r$ is the \emph{rank} of $G$, the dimension of $G/pG$ over $\mathbb F_p$. Homomorphisms $R^n\to G$ are given by the images of the basis vectors, so there are $|G|^n$ of them.

\begin{lemma}[counting surjections]\label{lem:countsur}
$\#\mathrm{Sur}(R^n,G)=|G|^n\prod_{i=0}^{r-1}(1-p^{i-n})$.
\end{lemma}

\begin{proof}
We show that $x_1,\dots,x_n\in G$ generate $G$ if and only if their images generate $G/pG$. One direction is clear. Conversely, if $H=\langle x_1,\dots,x_n\rangle$ satisfies $H+pG=G$, then the quotient $Q=G/H$ satisfies $pQ=Q$, hence $Q=p^kQ=0$. Now $G\to G/pG\cong\mathbb F_p^r$ is a surjective homomorphism, each fibre has $|pG|$ elements, so uniform independent $x_j$ have uniform independent images in $\mathbb F_p^r$. The images span $\mathbb F_p^r$ exactly when the $r\times n$ matrix they form has rank $r$, that is, has independent rows; choosing the rows one at a time gives $\prod_{i=0}^{r-1}(p^n-p^i)$ such matrices among $p^{rn}$. Hence the proportion of generating $n$-tuples is $\prod_{i<r}(1-p^{i-n})$.
\end{proof}

\begin{theorem}[moments of symmetric cokernels]\label{thm:symmoments}
Let $M$ be a uniform random symmetric $n\times n$ matrix over $R=\Z/p^k$, and $G$ as above. Then
\[
\mathbb E\,\#\mathrm{Sur}(\operatorname{cok}M,G)=\prod_{i=0}^{r-1}(1-p^{i-n})\cdot p^{\sum_{i<j}\min(\lambda_i,\lambda_j)}
\ \xrightarrow[n\to\infty]{}\ p^{\sum_{i<j}\min(\lambda_i,\lambda_j)}.
\]
\end{theorem}

\begin{proof}
One has $\mathbb E\,\#\mathrm{Sur}(\operatorname{cok}M,G)=\sum_{F}P(FM=0)$ over surjections $F:R^n\to G$, where $FM$ is the map $R^n\to G$, $y\mapsto F(My)$. We show that $P(FM=0)=p^{-n\sum\lambda_i+\sum_{i<j}\min(\lambda_i,\lambda_j)}$ for every surjective $F$; since $|G|^n=p^{n\sum\lambda_i}$, Lemma~\ref{lem:countsur} then gives the formula.

\emph{Step 1: every surjection is standard after a change of basis.} Let $\pi:R^n\to G$ send $e_i$ to $g_i$ for $i\le r$ and $e_i$ to $0$ for $i>r$. Given $F$ surjective, choose $x_1,\dots,x_r\in R^n$ with $F(x_i)=g_i$. Their images in $\mathbb F_p^n=R^n/pR^n$ are independent, because $F$ induces $\mathbb F_p^n\to G/pG$ sending them to a basis. Complete them to a basis of $\mathbb F_p^n$ and lift the added vectors to $x_{r+1},\dots,x_n\in R^n$. The matrix $A$ with columns $x_1,\dots,x_n$ has a determinant that is non-zero modulo $p$, hence a unit of $R$, so $A\in GL_n(R)$. Finally, for $j>r$ write $F(x_j)=\sum_ic_{ji}g_i$ with $c_{ji}\in R$ and replace $x_j$ by $x_j-\sum_ic_{ji}x_i$: this is a column operation, $A$ stays invertible, and now $F(x_j)=0$. So $FA=\pi$.

\emph{Step 2: the symmetric matrix does not notice.} Put $M'=A^{-1}M(A^{-1})^{T}$. The map $M\mapsto A^{-1}M(A^{-1})^T$ is a bijection of the set of symmetric matrices (its inverse is $M'\mapsto AM'A^T$), so $M'$ is again uniform. And $FM=\pi A^{-1}M=\pi M'A^T$, which vanishes if and only if $\pi M'=0$, since $A^T$ is invertible.

\emph{Step 3: count.} $\pi M'=0$ means $\sum_{i\le r}M'_{ij}g_i=0$ for every column $j$, that is, $p^{\lambda_i}\mid M'_{ij}$ for all $i\le r$ and all $j$. The independent entries of $M'$ are those with $i\le j$, each uniform on $R$, and a uniform element of $R$ is divisible by $p^a$ ($a\le k$) with probability $p^{-a}$. The conditions are: for $i\le r<j$, $p^{\lambda_i}\mid M'_{ij}$; for $i\le j\le r$, $p^{\max(\lambda_i,\lambda_j)}\mid M'_{ij}$ (the conditions from row $i$ and from row $j$ on the same entry); no condition when $r<i\le j$. So $P(\pi M'=0)=p^{-E}$ with
\[
E=(n-r)\sum_i\lambda_i+\sum_{i\le r}\lambda_i+\sum_{i<j\le r}\max(\lambda_i,\lambda_j).
\]
Using $\max(a,b)=a+b-\min(a,b)$ and $\sum_{i<j}(\lambda_i+\lambda_j)=(r-1)\sum\lambda_i$, we get $E=n\sum\lambda_i-\sum_{i<j}\min(\lambda_i,\lambda_j)$. The limit holds because the finite product tends to $1$.
\end{proof}

\begin{remark}[the exterior square]\label{rem:wedge}
The number $p^{\sum_{i<j}\min(\lambda_i,\lambda_j)}$ is the order of $\wedge^2G$, the quotient of $G\otimes G$ by the subgroup generated by the $x\otimes x$: indeed $\wedge^2(A\oplus B)\cong\wedge^2A\oplus\wedge^2B\oplus(A\otimes B)$, $\wedge^2$ of a cyclic group is $0$, and $\Z/p^a\otimes\Z/p^b\cong\Z/p^{\min(a,b)}$. This is how the theorem is usually stated: the moments of symmetric cokernels are $|\wedge^2G|$, where those of the cokernels of matrices with independent entries are $1$. The exterior square appears because a symmetric matrix defines a symmetric pairing on the torsion of its cokernel, and the conditions $FM=0$ only see the antisymmetric part of $G\otimes G$ once. For $G=\Z/p$ the moment is $1$, for $G=(\Z/p)^2$ it is $p$.
\end{remark}
